\documentclass[11pt,a4paper]{article}

\usepackage[T1]{fontenc}
\usepackage[utf8]{inputenc}
\usepackage[french]{babel}
\usepackage{geometry}
\usepackage{enumitem}
\usepackage{hyperref}
\usepackage{xcolor}
\usepackage{listings}

\geometry{margin=2.5cm}

\hypersetup{
    colorlinks=true,
    linkcolor=blue,
    urlcolor=blue
}

\lstset{
    basicstyle=\ttfamily\small,
    breaklines=true,
    frame=single,
    columns=fullflexible
}

\title{Projet --- Analyse de la Polarisation\\\large Point d'avancement technique}
\author{}
\date{\today}

\begin{document}

\maketitle

\section{Objet du document}

Ce document synth\'etise l'\'etat d'avancement du projet \emph{Analyse de la Polarisation} \`a ce stade. Il ne s'agit pas encore du rapport final demand\'e dans le sujet, mais d'un compte rendu de travail d\'ej\`a effectu\'e dans le d\'ep\^ot.

\section{Rappel du cadre du projet}

Le sujet demande de travailler sur deux types de profils \'electoraux :
\begin{itemize}[leftmargin=*]
    \item les votes par approbation, repr\'esent\'es par des vecteurs binaires ;
    \item les votes par ordres totaux, repr\'esent\'es par des permutations des candidates.
\end{itemize}

Le projet doit ensuite :
\begin{itemize}[leftmargin=*]
    \item g\'en\'erer des profils avec un niveau de polarisation contr\^olable ;
    \item impl\'ementer plusieurs mesures de polarisation ;
    \item impl\'ementer des distances entre bulletins ;
    \item proposer une estimation de la polarisation via un sch\'ema de type consensus global / consensus \`a deux groupes ;
    \item produire des exp\'eriences et des graphiques ;
    \item fournir un rapport final document\'e.
\end{itemize}

\section{Travail d\'ej\`a r\'ealis\'e}

\subsection{Structure du projet}

Une premi\`ere structure de code a \'et\'e mise en place pour s\'eparer les diff\'erentes parties du projet. Les fichiers suivants ont \'et\'e cr\'e\'es :

\begin{itemize}[leftmargin=*]
    \item \texttt{src/types.py}
    \item \texttt{src/generation.py}
    \item \texttt{src/distances.py}
    \item \texttt{src/pairwise.py}
    \item \texttt{src/phi2.py}
    \item \texttt{src/consensus.py}
    \item \texttt{src/clustering.py}
    \item \texttt{src/measures.py}
    \item \texttt{src/experiments.py}
    \item \texttt{src/plots.py}
    \item \texttt{scripts/run\_phi2\_experiments.py}
    \item \texttt{scripts/run\_distance\_experiments.py}
    \item \texttt{main.py}
    \item \texttt{README.md}
\end{itemize}

L'objectif de cette organisation est de faciliter la suite du projet : chaque partie du sujet peut \^etre reli\'ee \`a un module pr\'ecis.

\subsection{G\'en\'eration de profils}

Une premi\`ere version des g\'en\'erateurs al\'eatoires a \'et\'e impl\'ement\'ee dans \texttt{src/generation.py}.

\paragraph{Votes par approbation}
La fonction \texttt{generate\_approval\_profile} :
\begin{itemize}[leftmargin=*]
    \item g\'en\`ere un bulletin central binaire ;
    \item construit un bulletin oppos\'e bit \`a bit ;
    \item contr\^ole le niveau de polarisation \`a l'aide d'un param\`etre \texttt{alpha} ;
    \item ajoute du bruit avec un param\`etre \texttt{noise}.
\end{itemize}

L'id\'ee actuelle est la suivante :
\begin{itemize}[leftmargin=*]
    \item si \texttt{alpha} est faible, la majorit\'e des votantes restent proches d'un centre commun ;
    \item si \texttt{alpha} est forte, une part importante du profil est r\'epartie entre un centre et son oppos\'e.
\end{itemize}

\paragraph{Votes par ordres totaux}
La fonction \texttt{generate\_ranking\_profile} :
\begin{itemize}[leftmargin=*]
    \item g\'en\`ere un ordre central al\'eatoire ;
    \item utilise l'ordre inverse comme p\^ole oppos\'e ;
    \item contr\^ole le niveau de polarisation avec \texttt{alpha} ;
    \item ajoute du bruit par une suite d'\'echanges al\'eatoires.
\end{itemize}

Cette partie permet d\'ej\`a de pr\'eparer les exp\'eriences demand\'ees aux questions 1, 2, 6 et 15.

\subsection{Distances}

Une premi\`ere impl\'ementation des distances demand\'ees a \'et\'e ajout\'ee dans \texttt{src/distances.py}.

\paragraph{Distance de Hamming}
La fonction \texttt{hamming\_distance(a, b)} calcule le nombre de composantes diff\'erentes entre deux bulletins d'approbation de m\^eme taille.

\paragraph{Distance de Spearman}
Deux fonctions ont \'et\'e introduites :
\begin{itemize}[leftmargin=*]
    \item \texttt{ranking\_to\_positions(order)} convertit une permutation en tableau de positions ;
    \item \texttt{spearman\_distance(order1, order2)} calcule la somme des \'ecarts absolus de rang.
\end{itemize}

Ces fonctions constituent la base de travail pour les questions 7 et 8.

\subsection{Premi\`ere impl\'ementation de \texorpdfstring{$\phi_2$}{phi2}}

Une premi\`ere version de la mesure \(\phi_2\) a \'et\'e cod\'ee dans \texttt{src/phi2.py}, en s'appuyant sur \texttt{src/pairwise.py}.

\paragraph{Remarque importante}
Cette partie est encore consid\'er\'ee comme une \textbf{premi\`ere approximation de travail}. Elle permet d'avoir une cha\^\i ne ex\'ecutable pour les premiers tests et graphiques, mais la d\'efinition exacte de \(d^{c_k,c_l}(p)\) devra encore \^etre v\'erifi\'ee ligne \`a ligne par rapport au PDF avant de figer la version du rapport final.

Autrement dit :
\begin{itemize}[leftmargin=*]
    \item le squelette algorithmique est en place ;
    \item l'interpr\'etation math\'ematique fine de cette partie doit encore \^etre stabilis\'ee.
\end{itemize}

\subsection{Consensus global \texorpdfstring{$u_1$}{u1}}

Une premi\`ere version du calcul du consensus global a \'et\'e ajout\'ee dans \texttt{src/consensus.py}.

\paragraph{Cas approbation}
La fonction \texttt{approval\_consensus\_ballot} choisit pour chaque candidate la valeur majoritaire dans le profil. La quantit\'e \texttt{u1\_approval(profile)} somme ensuite les distances de Hamming entre ce consensus et tous les bulletins du profil.

\paragraph{Cas ordres totaux}
Une premi\`ere approximation de consensus est construite par rang moyen :
\begin{itemize}[leftmargin=*]
    \item on calcule pour chaque candidate son rang moyen dans le profil ;
    \item on trie ensuite les candidates selon ce rang moyen.
\end{itemize}

La fonction \texttt{u1\_ranking(profile)} somme alors les distances de Spearman entre ce consensus et les bulletins du profil.

\paragraph{Remarque importante}
Pour les votes par ordres totaux, cette version est une \textbf{version pragmatique initiale}. Le sujet indique qu'une approche plus rigoureuse doit s'appuyer sur un probl\`eme de couplage parfait de poids minimum. Cette partie devra donc \^etre raffin\'ee pour la version finale r\'epondant strictement \`a la question 11.

\subsection{Clustering \`a deux groupes et estimation de \texorpdfstring{$u_2$}{u2}}

Le fichier \texttt{src/clustering.py} contient une premi\`ere adaptation de l'algorithme de type \(k\)-means avec \(k=2\).

Le sch\'ema utilis\'e est :
\begin{enumerate}[leftmargin=*]
    \item initialisation al\'eatoire de deux centro\"\i des ;
    \item affectation de chaque bulletin au centre le plus proche ;
    \item recalcul du centro\"\i de de chaque groupe via le consensus du groupe ;
    \item r\'ep\'etition jusqu'\`a stabilisation ;
    \item r\'ep\'etition de plusieurs lancements et conservation du meilleur co\^ut trouv\'e.
\end{enumerate}

Deux fonctions ont \'et\'e ajout\'ees :
\begin{itemize}[leftmargin=*]
    \item \texttt{kmeans2\_approval}
    \item \texttt{kmeans2\_ranking}
\end{itemize}

Le r\'esultat retourn\'e contient actuellement :
\begin{itemize}[leftmargin=*]
    \item le co\^ut total ;
    \item les deux centro\"\i des ;
    \item les labels d'affectation.
\end{itemize}

\subsection{Mesures bas\'ees sur les distances}

Le fichier \texttt{src/measures.py} contient une premi\`ere impl\'ementation des mesures :
\begin{itemize}[leftmargin=*]
    \item \(\phi_{d_H}(p)\) pour les votes par approbation ;
    \item \(\phi_{d_S}(p)\) pour les votes par ordres totaux.
\end{itemize}

Les formules cod\'ees suivent l'id\'ee :
\[
\phi_{d_H}(p) = \frac{2}{nm}\bigl(u_1(p) - \widetilde{u_2}(p)\bigr)
\]
et
\[
\phi_{d_S}(p) = \frac{4}{nm^2}\bigl(u_1(p) - \widetilde{u_2}(p)\bigr).
\]

Ici, \(\widetilde{u_2}(p)\) est l'estimation obtenue par la proc\'edure it\'erative de clustering.

\subsection{Scripts d'exp\'eriences}

Deux scripts ont \'et\'e ajout\'es pour automatiser les premiers essais :
\begin{itemize}[leftmargin=*]
    \item \texttt{scripts/run\_phi2\_experiments.py}
    \item \texttt{scripts/run\_distance\_experiments.py}
\end{itemize}

Ils permettent :
\begin{itemize}[leftmargin=*]
    \item de faire varier le param\`etre \texttt{alpha} ;
    \item de r\'ep\'eter plusieurs tirages al\'eatoires ;
    \item de calculer moyenne et \'ecart-type ;
    \item de sauvegarder des fichiers CSV ;
    \item de produire des graphiques.
\end{itemize}

\subsection{Visualisation}

Le fichier \texttt{src/plots.py} contient une fonction \texttt{plot\_curve} qui trace une courbe avec barres d'erreur \`a l'aide de \texttt{matplotlib}. Cela pr\'epare directement les figures des questions 6 et 15.

\subsection{Point d'entr\'ee rapide}

Le fichier \texttt{main.py} fournit un test d'ex\'ecution minimal :
\begin{itemize}[leftmargin=*]
    \item g\'en\'eration d'un profil approval ;
    \item g\'en\'eration d'un profil ranking ;
    \item calcul de \(\phi_2\), \(\phi_{d_H}\) et \(\phi_{d_S}\) ;
    \item affichage de valeurs num\'eriques.
\end{itemize}

Ce point d'entr\'ee a \'et\'e ex\'ecut\'e avec succ\`es dans l'environnement courant.

\section{Tests et v\'erifications effectu\'ees}

\subsection{Installation des d\'ependances}

Le fichier \texttt{requirements.txt} a \'et\'e ajout\'e. Le paquet \texttt{pytest} a ensuite \'et\'e install\'e dans l'environnement courant afin de pouvoir lancer les tests automatiques.

\subsection{Jeu minimal de tests}

Les fichiers suivants ont \'et\'e ajout\'es dans \texttt{tests/} :
\begin{itemize}[leftmargin=*]
    \item \texttt{test\_generation.py}
    \item \texttt{test\_distances.py}
    \item \texttt{test\_phi2.py}
    \item \texttt{test\_consensus.py}
\end{itemize}

Ces tests v\'erifient actuellement :
\begin{itemize}[leftmargin=*]
    \item la forme des profils g\'en\'er\'es ;
    \item un cas simple pour la distance de Hamming ;
    \item l'annulation de la distance de Spearman sur deux ordres identiques ;
    \item l'annulation de \(\phi_2\) sur un profil parfaitement homog\`ene ;
    \item la coh\'erence de base entre \(u_1\) et l'estimation \`a deux groupes.
\end{itemize}

\subsection{R\'esultat de la campagne de tests}

La commande suivante a \'et\'e ex\'ecut\'ee :

\begin{lstlisting}
python -m pytest
\end{lstlisting}

R\'esultat observ\'e :
\begin{itemize}[leftmargin=*]
    \item 7 tests collect\'es ;
    \item 7 tests pass\'es.
\end{itemize}

\section{Etat actuel : ce qui est solide et ce qui reste \`a fiabiliser}

\subsection{Parties d\'ej\`a exploitables}

Les \'el\'ements suivants sont d\'ej\`a exploitables pour continuer le projet :
\begin{itemize}[leftmargin=*]
    \item la structure du d\'ep\^ot ;
    \item les g\'en\'erateurs de profils ;
    \item les distances de base ;
    \item une premi\`ere ex\'ecution de bout en bout ;
    \item des scripts d'exp\'eriences ;
    \item un socle de tests.
\end{itemize}

\subsection{Parties encore provisoires}

Les points suivants devront \^etre am\'elior\'es avant la r\'edaction finale :
\begin{itemize}[leftmargin=*]
    \item la d\'efinition exacte de \(d^{c_k,c_l}(p)\) et donc de \(\phi_2\) ;
    \item l'alignement strict de toutes les formules avec le PDF ;
    \item la partie th\'eorique des preuves pour les questions 4, 7, 9, 10 et 11 ;
    \item le calcul de \(u_1\) dans le cas des ordres totaux via une formulation conforme \`a l'indication du sujet ;
    \item la production des graphes finaux avec param\'etrage propre et commentaire r\'edactionnel.
\end{itemize}

\section{Prochaines \'etapes recommand\'ees}

L'ordre de travail recommand\'e pour la suite est le suivant :
\begin{enumerate}[leftmargin=*]
    \item corriger et stabiliser la d\'efinition math\'ematique de \(\phi_2\) ;
    \item aligner les fonctions \texttt{pairwise.py} et \texttt{phi2.py} avec le sujet ;
    \item raffiner le calcul de \(u_1\) pour les ordres totaux ;
    \item lancer les scripts d'exp\'eriences pour produire les premiers CSV et graphiques ;
    \item commencer la r\'edaction du vrai rapport final question par question.
\end{enumerate}

\section{Conclusion}

Le projet dispose maintenant d'une base de code propre, ex\'ecutable et test\'ee. Cette base ne constitue pas encore la version finale math\'ematiquement stabilis\'ee, mais elle permet de travailler efficacement sur la suite : mise au propre des d\'efinitions, exp\'eriences, graphiques et r\'edaction du rapport demand\'e.

\end{document}
